\documentclass[proposal,english]{AIGpaper}
% Please read the README.md file for additional information on the parameters and overall usage of AIGpaper

%%%% Package Imports %%%%%%%%%%%%%%%%%%%%%%%%%%%%%%%%%%%%%%%%%%%%%%%%%%%%%%%%%%%%%%%%%
\usepackage{graphicx}					           % enhanced support for graphics
\usepackage{amsfonts}					           % math fonts
\usepackage{amssymb}					           % math symbols
\usepackage{amsmath}					           % overall enhancements to math environment
\usepackage{amsthm}                                % for theorems and other environments

%%%% optional packages
\usepackage{url}                            % for linking external urls
\usepackage{booktabs}                      % for nicer tables 
\usepackage[hidelinks]{hyperref}           % for automatic hyperlinks
\usepackage{tikz}                          % creating graphs and other structures
\usepackage{todonotes} % for TODOs
\usepackage{csquotes} 


%%%% Enviroments for theorems, propositions, definitions, ....
\newtheorem{theorem}{Theorem}
\newtheorem{proposition}{Proposition}
\newtheorem{lemma}{Lemma}
\newtheorem{corollary}{Corollary}

\theoremstyle{definition}
\newtheorem{definition}{Definition}
\newtheorem{example}{Example}
\newtheorem{observation}{Observation}
\newtheorem*{observation*}{Observation}
%%%%%%%%%%%%%%%%%%%%%%%%%%%%%%%%%%%%%%%%%%%%%%%%%%%%%%%%%%%%%%%%%%%%%%%%%

%%%% Author and Title Information %%%%%%%%%%%%%%%%%%%%%%%%%%%%%%%%%%%%%%%%%%%%%%%%%%%
\author{Sebastian Frehmel}

%\title{Towards Live Streaming of Arbitrary 3D Gaussian Splatting Scenes: Velocity-Aware 3D Gaussians and Difference-Based Scene Updates in Fixed Multi-Camera Setups}
%\title{Velocity-Aware 3D Gaussian Splatting and Difference-Based Scene Updates as a Base for Live Streaming Arbitrary Scenes in Fixed Multi-Camera Setups}
%\title{Towards Live Streaming of 3D Gaussian Splatting: Velocity-Aware 3D Gaussians and Difference-Based Scene Updates in Fixed Multi-Camera Setups}

%\title{Towards Concurrent Real-Time\\ 3D Reconstruction and Novel-View Synthesis: \\Difference-Based Updates for 3D Gaussian Splatting Scenes in Fixed Multi-Camera Setups}
%\title{Difference-Based Updates for 3D Gaussian Splatting Scenes in Fixed Multi-Camera Setups \\\vspace{0.3em} {\small\textit{or alternatively}} \vspace{0.3em} \\Towards Unrestricted Real-Time\\ 3D Reconstruction of Dynamic Scenes: \\Difference-Based Updates for 3D Gaussian Splatting Scenes in Fixed Multi-Camera Setups}
%Isolating and Updating Active Regions\\ of Dynamic 3D Gaussian Splatting Scenes\\ in Fixed Multi-Camera Setups

\title{Incremental Subscene Updates for Unconstrained Online Dynamic Scene Reconstruction via 3D Gaussian Splatting in Fixed Multi-Camera Setups} %\\\vspace{0.3em} {\small\textit{or alternatively}} \vspace{0.3em} \\Online Dynamic Scene Reconstruction\\ with 3D Gaussian Splatting: \\Training Small Delta Scenes\\ in Fixed Multi-Camera Setups}

%\title{Real-Time, Difference-Based Updates\\ to 3D Gaussian Splatting Scenes\\ in Fixed Multi-Camera Setups}
%Real-Time Updates to 3D Gaussian Splatting Scenes in Fixed Multi-Camera Setups Using a Difference-Based Approach
%Real-Time 3D Gaussian Splatting Updates in Fixed Multi-Camera Setups Using a Difference-Based Approach
%Updating 3D Gaussian Splatting Scenes in Real-Time Using a Difference-Based Approach
%Updating 3D Gaussian Splatting Scenes in Fixed Multi-Camera Setups in Real-Time Using a Difference-Based Approach

%%%% Abstract %%%%%%%%%%%%%%%%%%%%%%%%%%%%%%%%%%%%%%%%%%%%%%%%%%%%%%%%%%%%%%%%%%%%%%

\englishabstract{Online dynamic scene reconstruction via 3D Gaussian Splatting has seen comparatively little research, and existing works typically impose strict constraints on scene geometry.
To address this limitation, the proposed thesis introduces a novel algorithm that incrementally updates the full 3D scene exclusively through small subscenes learned from frame-to-frame pixel changes. The focus lies on algorithmic feasibility and on performance evaluation in real-time contexts.}


\begin{document}

\maketitle % prints title and author information, as well as the abstract 


% ===================== Beginning of the actual text section =====================
\section{Introduction}
This chapter showcases the various relevant technologies for the proposed thesis.

\subsection{Brief Overview of 3D Gaussian Splatting (3DGS)}
3DGS is a young technology in the field of 3D Reconstruction which is itself part of the wider Computer Graphics field. It was first published by Kerbl et al.\ at INRIA in France in 2023 \cite{kerbl3DGaussianSplatting2023} and has since experienced wide adoption and received critical acclaim across domains, having been cited more than 10,000 times \cite{semanticScholarKerbl,researchGateKerbl}. 3DGS solved issues previous technologies in 3D Reconstruction had, such as Neural Radiance Fields (NeRFs), as for example being able to render 1080p resolution scenes in real-time (\textit{Novel-View Synthesis}), making full use of modern GPGPUs, offering editability of the scene while still featuring very reasonable training times. 3DGS uses 3D Gaussians as the base primitive, a previously known data structure, which can be described as ellipsoid fuzzy blobs that can take any ellipsoid shape from perfect spheres to thin and long shapes or flat and round shapes. They use a multi-dimensional normal distribution for their opacity distribution, being most opaque at the center and rapidly losing opacity towards the edges. Kerbl et al., among other achievements, made 3D Gaussians differentiable so they can be used in standard backpropagation with gradient descent. This makes 3DGS an explicit, learnable radiance field representation without neural networks.

\subsection{3D Gaussian Splatting Algorithm}
Regular 3DGS takes 2D input images and learns the pictured 3D scene from them. The training phase can run for 2-20 minutes, depending on scene size but already good results are available within seconds after the start of training. A large 3D scene can have millions of 3D Gaussians and hundreds of input images. The input images are first converted into a point cloud and camera matrices (camera orientation in space etc.)\ using an algorithm called \textit{Structure from Motion} \cite{ullmanInterpretationStructureMotion1979} and the points of the point cloud become the origins for the initial 3D Gaussians before training. It is important to know that even initialising 3D Gaussians randomly still converges, but takes longer to train. Each forward and backward pass processes one random input image to adjust the 3D scene iteration by iteration until the loss between original input images and rendered images is small enough. The backward pass adjusts all attributes of a 3D Gaussian (position, opacity, shape, color) purely by the difference in pixel color between rendered and original image after the previous froward pass.

I attended the AIG seminar \enquote{Artificial Intelligence - Visual Computing} with 3DGS as the topic, covering the base concept and algorithm in much more detail.

\subsection{Dynamic Scene Reconstruction}
Dynamic scene reconstruction (as in \enquote{reconstructing a dynamic scene}, not as in \enquote{dynamically reconstructing a scene}) in general represents reconstruction of 3D scenes that change over time, which means the input is a video feed with a certain frame rate. 

\subsubsection{Offline Dynamic Scene Reconstruction}
In the context of 3DGS as a technology, dynamic scene reconstruction is also called \enquote{4DGS}. Training the initial scene is identical to regular 3DGS, but the following frames can rely on scene geometry from the previous frame, making training a lot faster as 3D Gaussians are already in the vicinity, with similar shape and color of the next frame's content. It can still not be done in real-time, though, as the input imagery is still full 2D images and the full 3D scene is being considered for training. The difference of the 3D scene between frames can be encoded into/learned by so-called \textit{deformation fields}, small neural networks which allow transforming 3D Gaussians over time according to the field. The other option is to encode the time component directly into the 3D Gaussian data structure, effectively transforming them into 4D Gaussians. This means that 4DGS scenes can only be learned from pre-recorded imagery, then encoded and then later replayed and explored with novel-view synthesis. There is no real-time component during reconstruction\footnote{Novel-view synthesis remains real-time, as with regular 3DGS}, which is why the attribute \enquote{offline} can be assigned to this scenario.

\subsubsection{Online and Real-Time Dynamic Scene Reconstruction}
Contrary to 4DGS scenes, online dynamic scene reconstruction (also called \enquote{3DGS Live Streaming}, \enquote{Live 3DGS}, \enquote{Live 4DGS} or \enquote{3DGS Live Video}) describes concurrent reconstruction and novel-view synthesis, usually after initial regular training of a base 3DGS scene. Online dynamic scene reconstruction does not have to be real-time, but it can be, which means that real-time dynamic scene reconstruction represents the strictest class in dynamic scene reconstruction.

In both cases, a set of new input imagery is converted into the next iteration of the full 3D scene using various algorithms on the fly. There is no need for deformation fields or 4D Gaussians, even though they have been used to store the dynamic scene for later.

Real-time dynamic scene reconstruction has so far not been achieved for unconstrained/\allowbreak arbitrary scenes and has seen some solutions imposing major constraints on the allowed geometry. Existing solutions e.g.\,use AI priors (large pre-trained neural networks) to predict the next movement in the scene, limiting what can appear in the scene, or use skeletal models of humans to guide 3D Gaussians.

\subsubsection{Camera Setups for Dynamic Scene Reconstruction}
In general, dynamic scene reconstruction scenarios can be divided into single or multi-camera setups. Single cameras can be found in 3DGS solutions to SLAM (Simultaneous Localization and Mapping, mobile camera) in robotics and live 3D telepresence (fixed camera) for example. 

For the proposed thesis a fixed multi-camera setup is used in order to be able to compare consecutive frames. Mobile cameras add much to the general computational load as their position has to be estimated with every frame, which is why using fixed multi-camera setups can help bound the load. This automatically limits the extent of the scene but comes naturally for e.g.\,confined indoor spaces.

Another option within this differentiation is so-called \textit{event-based 3DGS} which uses \enquote{event cameras} that only record changes in light intensity, without color. State-of-the-art solutions for SLAM or robotics in general use these kind of cameras allowing real-time 3D reconstruction of unconstrained scenes, however, at lower resolutions and without color. The proposed thesis takes a different approach but solutions in SLAM, especially with event-based 3DGS, could prove to be valuable input, too.

Another option fascilitating dynamic reconstruction are RGB-D cameras, which record depth information alongside color values. This however works only for small scenes as larger scenes require Lidar-based sensors, which have their own restrictions and side-effects.

\subsection{Overview of this Proposal}
%\todo[inline]{Write here an introduction for the topic of your planned thesis. Describe why is considering/investigating the topics of the thesis important.}

The proposed thesis is about formulating and evaluating a novel approach to online dynamic scene reconstruction via 3DGS that only looks at the subset of 3D Gaussians - the \textit{subscene} - that is being altered in comparison to the previous frame. Such a subscene would contain only a fraction of all Gaussians and would therefore be much faster to compute. At first, though, the space covered by the fixed multi-camera setup is trained as a regular 3DGS scene which functions as the base for subsequent updates.

The input imagery for a subscene are frame-to-frame pixel changes from the new frame to the previous frame (for each camera). This is explicitly not a difference image, but the input image would contain pixels that did change over two consecutive frames holding their changed RGB value. The rest of the pixels would be encoded as 'unset' and be ignored during subscene reconstruction.

During the forward pass of subscene reconstruction, only the set pixels (the ones not to be ignored) are considered and the loss function only compares the set pixels in rendered and input image before heading into the backward pass. Also in the backward pass, only Gaussians behind set pixels are considered - this is provided by 3DGS out of the box. All other Gaussians would not be touched during gradient descent (\enquote{frozen}).

With the existing geometry being close to what the new imagery of a new frame requires, the algorithm is expected to converge very quickly. New geometry appearing where there was none in the scene before is most likely very small in the first frame it appears in, so it can just be reconstructed directly. Geometry being removed/leaving the scene will be handled in the regular way as the pixel color will then show the color behind the geometry that was removed. This geometry will have been known before and the algorithm should only remove the 3D Gaussians that were in front of the one that provided the color behind.\\
This means that the training algorithm for the subscene is essentially the same as with regular 3DGS, \enquote{only} using a mask on pixels and Gaussians to be considered. This process could be sped up if Gaussians contained motion vectors making predicting the Gaussians' next positions even more precise.

This subscene approach in a fixed multi-camera setup has never been investigated but looks very promising to at least reach a low number of FPS which can then later -- in subsequent work -- be refined to run at real-time speeds (\enquote{First make it work, then make it work fast}). 3DGS in general is a 'hot topic' in academic research and the reason for the subscene approach not having been considered so far in dynamic scene reconstruction is at least two-fold: single-camera setups, especially for robotics/SLAM are a major focus of research and researchers are very focused on using other AI/ML methods for scene prediction, as it brings grant money more easily and, too, looks promising.



%So this proposal is essentially about improving 4DGS to work under real-time constraints. This means we have a static 3D scene that is first being learned using the unmodified, regular 3DGS algorithm\footnote{Which iteration of the algorithm has still to be decided}. Here we can see why it is important to work with fixed multi-camera setups in real-time settings. With fixed cameras, the camera's field of view can't change, the camera matrix is always the same and it's thus simple to create differences between two subsequent captured images.



%For the real-time component we can deduce that we get 33 milliseconds per frame to advance the scene if we consider 30 frames per second as smooth enough. For a decently sized static initial scene that takes 5 minutes to train, each iteration of the training algorithm (there are usually around 30,000 iterations) takes 10ms. So the goal is to fill the available 33ms with efficient (GPU) logic to fully advance the scene by one frame. Seeing that the full scene takes 10ms for one training iteration the consequence is clear: what is being considered for the real-time calculation has to be much less than the full scene in order to work at all. However, in 33ms, a modern Nvidia consumer GPU like the GeForce 5090 can perform $3.45$ \textit{trillion} FP16 operations fully in parallel while datacenter GPUs like the Nvidia B200 can even reach 74.5 trillion operations (this is the theoretical maximum, 40-60\% of this are usually considered \enquote{good}).

%This is where the novel approach steps in. For the dynamic part of the 3D scene (the one after training the initial, static 3D scene) I propose to only use the differences between subsequent captured images of each camera to contribute to the scene. Fixed cameras will only ever see the same section of the scene, allowing only the changes (the \enquote{diff}) in the scene to propagate into the novel algorithm.

%We also store the previous diff and as a first step, those two diffs are compared and similar/same pixes are mapped onto each other.
%What the 3DGS algorithm does out of the box is map pixels to 3D Gaussians (1:n). So when we have a pixel, we know which 3D Gaussians will be affected. We mark all affected 3D Gaussians as \textit{active} and separate them from the full scene (+ some neighboring 3D Gaussians as padding
% and 3D Gaussians in the direction the pixel was moving
%). Then the centerpiece of the novel algorithm iterates over the diff-based input images to create a small, possibly tiny, new 3D scene, which will generate quickly as the active 3D Gaussians represent shapes and positions very close to the ones in the new frames. This small diff-based 3D scene is then inserted back into the full scene. As we only touched affected 3D Gaussians, this should not result in gaps etc.


%Geometry being removed/leaving the scene will be handled in the regular way as the pixel color will then show the color behind the geometry that was removed. This geometry will have been known before and the algorithm should only remove the 3D Gaussians that were in front of the one that provided the color behind.

%learning directly from encoded sparse image?

%RLE?


\section{Background}
\todo[inline]{In this section, describe the background for the topic of your thesis. Generally, make use of the features of documents (and LaTeX) by providing figures, tables, referring, and other abilities to structure the text.}

20 minutes = 1200000 ms, divided by 30000 iterations = 40ms per iteration

5 minutes = 10ms per iteration

Look into SLAM applications

Für weitere Informationen zum Darstellen von Argumentationsgraphen, siehe\footnote{\tiny\ttfamily\url{https://github.com/aig-hagen/tikz_argumentation/blob/main/argumentation-doc.pdf}}.

% \begin{table*}[t]
% \begin{center}
% \begin{tabular}{lll}
% \hline
% 1&2&3\\
% 4&5&6\\
% \hline
% \end{tabular}
% \end{center}
% \caption{Eine Tabelle.} \label{tab:t1}
% \end{table*}


\subsection{Encoding Test}
Städte, Länder, Flüsse


\section{Problem Description}
\todo[inline,caption={}]{Describe what the premises and problems in the thesis are dealing with. 
	Generally, this is often some lack of knowledge or lack of means that is planned to be resolved. For instance, non-existing insights in a certain domain or no tools/algorithms/implementations for further investigations.
}
For propositions and definitions please use appropriate environments.
Here is an example for a definition.

\begin{definition}
	A deterministic finite automaton (DFA) is a tuple \( A=(Q,\Sigma,\delta,q_0,F) \) consisting of
	\begin{itemize}
		\item a finite set of states \( Q \),
		\item a finite alphabet \( \Sigma \),
		\item a transition function \( \delta:Q\times\Sigma\to Q \), 
		\item a start state \( q_0 \in Q \), and
		\item a set of accepting states \( F \subseteq Q \).
	\end{itemize}
\end{definition}

Here is an example for a proposition.

\begin{proposition}
	The problem of deciding whether a Turing machine halts for an input string is undecidable.
\end{proposition}

Examples can be made with the example environment.

\begin{example}
	Content of the example environment.
\end{example}


\section{Goals of the thesis}
\todo[inline,caption={}]{In this section, describe what are the plans and goals for the thesis. Contributions are existing in many forms. Generally, as time for thesis is quite short, we recommend to be conservative with settings goals.
	Contributions which are part of \textbf{every} thesis are:
	\begin{itemize}
		\item Literature research and description of the background in a proper and concise way
	\end{itemize}
Here are some examples for other types of contributions that can be also goals of a thesis:
\begin{itemize}
	\item Identification and documentation of problems
	\item Illustration of definitions, theorems and algorithms
	\item Development, identification and presentation of examples
	\item Evaluations in different forms (empirical and theoretical)
	\item Implementations
	\item Development or investigations of concepts
\end{itemize}
}

\section{About me}
I am 43, from Berlin and I studied Diplom-Informatik at Karlsruhe Institute of Technology, graduating in 2010. After that I spent around twelve years working as a software developer, mostly on the data and backend sides of projects. I decided to return to enrolling at a university and studying Data Science mostly because of two long-term diseases that prevented me from performing my previous jobs (including a PhD program) but also from the felt need to get up to speed on more recent topics in Computer Science. As an experienced software developer with a Diplom in CS, I focussed heavily on all modules offered by AIG as they offered the most added value compared to \enquote{regular} CS modules. Keen on computer visuals, I was happy to learn about 3D Gaussian Splatting (3DGS) in late 2025 and have since tried to focus on this area as a thesis topic. 3DGS offers the challenge I seek at the intersection of advanced Computer Science and ML/AI. The idea of the thesis topic stems from an older thought experiment of mine, which saw true 3D video broadcasts of live events become possible.\\
I attended the seminar \enquote{Artificial Intelligence - Visual Computing} in the 2026 summer term, actually proposing 3DGS as a seminar topic early on. The seminar paper I then wrote on 3DGS forms an important knowledge base of this proposal and the proposed master thesis. Without the seminar paper, such a master thesis would probably not have been feasible due to time constraints.
%If this thesis shows that my proposed approach is possible, I can think of turning the concepts into a business idea.

\newpage

\bibliographystyle{plain}
\bibliography{references}

\end{document}
